\subsection{Confidence Intervals}
\paragraph*{Main Idea:}
If $\bar{X}$ is close to $\mu$ with 95\% probability, then $\mu$ is also close to $\bar{X}$ with 95\% probability.
Using this idea, we can use our point estimate $\bar{X}$ to create a \textbf{confidence interval} for $\mu$.

\paragraph*{Definition:}
After observing a random sample from a normal population of size $n$ and computing the sample mean $\bar{X}$,
a 95\% confidence interval for $\mu$ itself
\[
    CI_\mu = (\bar{X} - 1.96 \cdot \frac{\sigma}{\sqrt{n}}, \bar{X} + 1.96 \cdot \frac{\sigma}{\sqrt{n}})
\]

This means that with 95\% probability, the interval $CI_\mu$ contains the true mean $\mu$.

\paragraph*{Example:}
A random sample of $n=50$ males showed an average daily intake of dairy products equal to 756 grams.
Suppose it is known that dairy intake follows a normal distribution with standard deviation of
35 grams. Fins a 95\% confidence interval for the population mean $\mu$.

\[
    CI_\mu = (\bar{X} - 1.96 \frac{\sigma}{\sqrt{n}}, \bar{X} + 1.96 \cdot \frac{\sigma}{\sqrt{n}}) = (756 - 1.96 \cdot \frac{35}{\sqrt{50}}, 756 + 1.96 \cdot \frac{35}{\sqrt{50}}) = (744.26, 767.74)
\]

\paragraph*{Interpreting Confidence Intervals:}
\textbf{Warning:} Assume that you collect data ($X_1, \dots, X_n$) and compute a 95\% confidence interval.
It is tempting, but wrong, to say that your confidence interval contains the true mean $\mu$ with 95\% probability.
Since the true mean is a fixed number, it will be contained in any interval with probability either 0 or 1. The
\textbf{Confidence Level} (here 95\%) rather refers to the generation of the random variables or the sampling process.

\paragraph*{Confidence Levels:}
To compute conficence intervals for levels other than 95\%, we have to change the critical z-value from 1.96 to 
another number. You can find the appropriate number in the Normal CDF table. Therefore, the general formula for a
confidence interval with confidence level $100(1 - \alpha)\%$ is:
\[
    CI_\mu = (\bar{X} - z_\frac{\alpha}{2}\frac{\sigma}{\sqrt{n}}, \bar{X} + z_\frac{\alpha}{2}\frac{\sigma}{\sqrt{n}})
\]

\paragraph*{Confidence Level, Precision, and Sample Size:}
Which is better? A 95\% confidence interval or a 99\% confidence interval? It depends on how much precision is required. 
Higher confidence levels lead to wider intervals, therefore less precision. If more precision is required, a larger sample size is needed.
The following factors influence the width of a confidence interval:
\begin{itemize}
    \item Confidence Level: Higher confidence levels lead to wider intervals.
    \item Sample size: Larger sample sizes lead to narrower intervals.
    \item Standard deviation of the population: Larger standard deviations lead to wider intervals.
\end{itemize}